\chapter{Overview}
\label{ch:overview}

\section{What a .clinerules agent is}

A \artifact{.clinerules} file is a local contract for an agent. It is
not application code and not a product requirement document. It is the
bridge between the two:
\begin{itemize}
  \item the specification states what must be true;
  \item the \artifact{.clinerules} file states how an agent must work
    in this repository to keep that true.
\end{itemize}

\section{Two mandatory agent types}

Every governed domain should have two mandatory agent contracts.

\subsection*{1. Development agent}

The development agent owns implementation, integration, migration, test
obligations, and completion criteria. Its contract lives in
\filepath{src/<domain>/.clinerules}.

\subsection*{2. Runtime-monitor agent}

The runtime-monitor agent owns health criteria, alert conditions,
evidence correlation, and escalation behaviour. Its contract lives in
\filepath{src/<domain>/.clinerules.runtime-monitor}.

\section{Why the split is required}

\begin{adrbox}{Separation of construction from supervision}
A single rule file is not enough for a governed agentic platform. The
development agent must optimize for implementation correctness and
contract coherence. The runtime-monitor agent must optimize for
observability, alerting precision, and fail-closed behaviour. Keeping
them separate avoids blending build-time convenience with runtime risk
control.
\end{adrbox}

\section{Why a swarm is needed}

The artefacts in \filepath{docs/pdf/specs/} are cross-domain. They span
Arabic, TB, regulations, Simula, shared fragments, and the
implementation rule packs that keep those books honest. A single agent
can draft changes, but a reliable delivery model needs multiple roles:
\begin{itemize}
  \item domain specialists who know their code and contracts;
  \item regulatory specialists who impose the governance envelope;
  \item runtime monitors that validate live control assumptions;
  \item documentation and PDF build specialists who assemble and
    validate the published specifications.
\end{itemize}

\section{Rule-pack lifecycle states}
\label{sec:lifecycle-states}

A rule-pack transitions through defined lifecycle states to ensure
governance integrity before enforcement begins.

\begin{table}[htbp]
\centering
\footnotesize
\caption{Rule-pack lifecycle states.}
\label{tab:lifecycle-states}
\begin{tabularx}{\textwidth}{@{}l l X@{}}
\toprule
\textbf{State} & \textbf{Trigger} & \textbf{Behaviour} \\
\midrule
\texttt{DRAFT} & Feature branch with \texttt{status: draft} frontmatter & Structural validity MAY be checked; CI/CD gates SHALL NOT block on draft packs. \\
\texttt{REVIEW} & Pull request opened against default branch & Full validation harness runs; reviewers must approve before merge. \\
\texttt{ACTIVE} & Merged to default branch & Pack is enforced; all agents must comply with its obligations. \\
\texttt{DEPRECATED} & \texttt{status: deprecated} frontmatter added & Pack remains enforced for one release cycle; agents receive deprecation warnings. \\
\texttt{RETIRED} & Removed from repository & Pack is no longer enforced; historical references remain in audit logs. \\
\bottomrule
\end{tabularx}
\end{table}

\subsection{Lifecycle invariants}

\begin{itemize}
  \item A rule-pack SHALL be considered \textbf{active} only when merged
    to the default branch without a \texttt{status: draft} or
    \texttt{status: deprecated} frontmatter field.
  \item During development, a rule-pack MAY reside in a feature branch
    and be marked as \texttt{DRAFT} via frontmatter. Draft packs SHALL
    NOT block CI/CD gates, but their structural validity MAY be checked.
  \item Deprecation requires a minimum one-release window before
    retirement to allow dependent agents to adapt.
  \item Retirement does not erase audit history; all actions taken under
    the pack remain attributable.
\end{itemize}
